\chapter{Related Work}

The intersection of machine learning and financial text analysis has evolved substantially over the past two decades. Early approaches relied on classical bag-of-words representations and linear models, while more recent work leverages deep neural networks and transformer-based architectures. In parallel, advances in representation learning have enabled increasingly sophisticated modeling of long-form financial documents such as SEC filings. This chapter reviews the most relevant strands of machine learning research related to stock price prediction using textual and market data, with particular emphasis on deep learning and domain-adapted language models.

One of the foundational works in text-based stock prediction by Tetlock (2007)\cite{tetlock2007giving} examined whether the tone of media content predicts market movements. Although not a deep learning paper, this work demonstrated that linguistic sentiment extracted from textual data contains economically meaningful signals. Textual tone was measured using dictionary-based methods, and regressions were used to link media pessimism to market returns. The key contribution was to empirically establish that textual sentiment correlates with short-term market behavior.

Subsequent research began incorporating supervised machine learning techniques. Schumaker et al. (2009)\cite{schumaker2009textual} applied support vector machines and naive Bayes classifiers to financial news articles, demonstrating that textual features could improve short-term stock direction prediction. Their objective was to classify whether a stock would increase or decrease following news releases. Using engineered linguistic features and classification algorithms, they reported predictive accuracy exceeding naive baselines, suggesting that machine learning could capture nonlinear relationships between text and returns.

These early works relied heavily on manually curated lexicons or shallow feature representations such as TF-IDF vectors. While effective to a certain degree, such representations fail to capture contextual semantics and long-range dependencies present in financial reports.


The introduction of deep learning marked a turning point in financial text modeling. Ding et al. (2015)\cite{ding2015deep} proposed a deep neural network framework that combined event embeddings derived from news headlines with temporal price modeling. Their model used a convolutional neural network (CNN) to encode event representations and incorporated historical price information. The objective was to predict stock price movements following news events. They demonstrated that deep representations significantly outperformed traditional bag-of-words approaches.

Similarly, Hu et al. (2018)\cite{hu2018listening} introduced a hierarchical attention network for modeling financial news and predicting stock volatility and direction. By applying attention mechanisms, the model could weight more informative sentences more heavily, thereby improving interpretability and performance. This work highlighted the importance of document-level structure in financial text modeling.

Parallel to developments in text modeling, recurrent neural networks (RNNs), particularly long short-term memory (LSTM) networks, were widely applied to time-series financial forecasting. Fischer et al. (2018)\cite{fischer2018deep} demonstrated that LSTMs outperform traditional autoregressive models for stock return prediction using price-based features. Although their work focused primarily on numerical time-series data rather than textual inputs, it established deep learning as a competitive alternative to classical econometric methods.

These contributions collectively show that neural architectures can effectively capture nonlinear patterns in financial data. However, many of these models were applied to relatively short textual inputs, such as headlines or summaries, rather than full-length regulatory filings.

The development of transformer architectures (2017)\cite{vaswani2017attention} and the introduction of BERT (2019)\cite{devlin2019bertpretrainingdeepbidirectional} revolutionized natural language processing. BERT's bidirectional self-attention mechanism enables contextual word representations, allowing models to capture nuanced semantics and long-range dependencies. This was particularly relevant for financial text, which is formal, domain-specific, and context-dependent.

Recognizing that general-purpose language models may not optimally capture financial terminology, researchers began developing domain-adapted transformer models. One of the most influential contributions is FinBERT (2020)\cite{yang2020finbert}, which fine-tunes BERT on financial corpora, including SEC filings and financial news. The objective of FinBERT was to improve sentiment classification accuracy in financial contexts. Empirical results showed that domain adaptation significantly improves performance compared to general BERT models when applied to financial sentiment tasks.

The existing literature demonstrates several key trends: (i) the transition from lexicon-based sentiment analysis to deep contextual embeddings, (ii) the growing importance of domain-specific language models, and (iii) the integration of textual and structured financial data within multimodal frameworks.

However, several limitations remain in prior work. First, many studies focus exclusively on news articles rather than regulatory filings, which are more comprehensive and standardized disclosures. Second, when SEC filings are used, models often truncate documents or rely on limited sections without systematic aggregation strategies. Third, relatively few works explicitly examine sector and industry-adjusted returns as alternative prediction targets, thereby limiting insight into whether models capture idiosyncratic or systematic components of stock movements.

The present study differs from prior research in several important ways. First, it focuses specifically on 10-K and 10-Q filings from U.S. markets over a recent and standardized period (2011--2024), ensuring consistency in reporting format. Second, rather than truncating long documents, it employs a controlled chunk-based sampling and aggregation strategy to approximate document-level sentiment using FinBERT. This approach balances computational feasibility with broad coverage of narrative disclosures. Third, it systematically evaluates multiple prediction horizons (1--7 days) and considers both regression and multi-class classification formulations. Finally, by constructing sector and industry-adjusted return labels using a leave-one-out methodology, this study explicitly distinguishes between firm-specific and systematic market effects.

Taken together, this work contributes to the literature by providing a comprehensive and methodologically rigorous framework for studying the short-term market impact of financial disclosures using modern transformer-based language models combined with structured price-based features.